\section{Baseline model evaluation}

\subsection*{Apply the MVG model to the project data. Compute the corresponding error rate.}
Running the Standard Multivariate Gaussian (MVG) model on the full 6-dimensional dataset yields an accuracy of \textbf{93.0\%}, which corresponds to an error rate of \textbf{7.0\%}. 

The standard MVG model achieves this strong performance because it computes a full, class-specific covariance matrix. This allows the model to capture both the individual variance of each feature and the covariances between them, independently for the "Fake" and "Genuine" classes. As observed in Laboratory 4, features 1, 2, 3, and 4 are reasonably well-approximated by a Gaussian distribution. Even though features 5 and 6 exhibit heavy skewness and sharp drop-offs, the unconstrained MVG is still able to extract enough discriminative information from the remaining features to achieve the lowest baseline error rate.

\subsection*{Apply now the tied Gaussian model, and compare the results with MVG and LDA. Which model seems to perform better?}
The Tied Covariance MVG model results in an error rate of \textbf{9.3\%}. 

Compared to the Standard MVG (7.0\%), the Tied model performs significantly worse. However, when we compare it to Linear Discriminant Analysis (LDA), the Tied MVG performs identically, as the LDA model also achieved an error rate of \textbf{9.3\%}. 

This identical performance makes perfect theoretical sense: both LDA and Tied MVG rely on the exact same fundamental mathematical assumption. They both assume that all classes share the same within-class covariance matrix, resulting in a strictly linear decision boundary. 

The Standard MVG performs better than both because the underlying assumption of a shared covariance is false for this dataset. As we discovered in our exploratory analyses (Laboratories 2 and 4), the "Fake" and "Genuine" classes exhibit significantly different variances for several features (for example, Genuine Feature 1 has a much higher variance than Fake Feature 1). By forcing a single, pooled covariance matrix, the Tied MVG and LDA destroy this crucial geometric distinction. The unconstrained MVG successfully exploits these variance differences using a quadratic decision boundary, making it the better-performing model.

\subsection*{Finally, test the Naive Bayes Gaussian model. How does it compare with the previous two?}
The Naive Bayes model achieves an error rate of \textbf{7.2\%}. 

It performs remarkably well—vastly outperforming the Tied MVG and LDA (9.3\%) and coming extremely close to the fully parameterized Standard MVG (7.0\%). 

The Naive Bayes model simplifies the standard MVG by assuming that all features are conditionally independent within each class, meaning it mathematically forces all out-of-diagonal elements in the covariance matrices to zero. The fact that the error rate only degrades by a negligible 0.2\% compared to the full MVG strongly implies that the features in this dataset are already largely uncorrelated. Because the cross-feature correlations do not carry significant discriminative weight, the Naive Bayes independence assumption holds up well in practice. It easily surpasses the Tied model because it retains the class-specific variances that the Tied model discards, proving that capturing distinct variances is much more important for this dataset than capturing cross-feature correlations.